\chapter*{Índice Comentado} 

La presente memoria se organiza en cuatro ejes principales. En la Parte I: Contextualización (página~\pageref{part:contexto}), donde es presentado el estado actual de la generación de energía eléctrica, los antecedentes relevantes en la materia y la problemática concreta que motiva el desarrollo de este trabajo. Esta parte concluye con la definición del objetivo general y los objetivos específicos que guían el resto del documento.

La Parte II: Investigación (página~\pageref{part:investigacion}) reúne los fundamentos necesarios para comprender el diseño del sistema, es decir, los conceptos teóricos, protocolos de comunicación y tecnologías de medición y de \ing{software} sobre los que se apoya la propuesta. A partir de estos fundamentos se describe la arquitectura general del sistema y sus componentes, y complementándose con un análisis del mercado en el que la solución propuesta en este trabajo podría insertarse.

La Parte III: Desarrollo y Resultados (página~\pageref{part:desarrollo}), describe la implementación del \ing{hardware} diseñado, el \ing{firmware} desarrollado y la plataforma de supervisión remota, así como la integración de todas estas partes y los ensayos experimentales realizados para validar el funcionamiento del sistema desarrollado en condiciones representativas de operación.

La Parte IV: Discusión y Conclusiones (página~\pageref{part:conclusiones}), analiza los resultados obtenidos frente a los objetivos planteados, discute las limitaciones encontradas durante el desarrollo y propone líneas de trabajo futuro. El documento concluye con los Anexos, que contienen la documentación técnica complementaria y la bibliografía consultada.